\section{Resultados por indicador}
\label{sec:resultados}

A Figura~\ref{fig:contrasteoss} resume o contraste central da pesquisa:
quanto cada indicador difere entre a gestão por OSS e a administração
Direta, já descontados complexidade, porte e ano, com as margens de
erro agrupadas por hospital.

\textbf{Uma advertência precisa acompanhar todos os números desta seção,
e ela não está nas margens de erro.} Os coeficientes estimados \emph{dentro}
do hospital (o tempo de permanência e o faturamento) são identificados por
\textbf{cinco conversores}, e a Subseção~\ref{sub:oqueforam} documenta que
nenhum dos cinco é um caso limpo: uma conversão foi decidida em resposta a
uma crise de gestão, outra coincidiu com mudança de prédio, dois hospitais
já operavam com serviços terceirizados antes da mudança formal, e os dois
únicos com janela posterior longa foram para a mesma organização social.
Os intervalos de confiança medem incerteza amostral e não sabem nada
disso.

O quadro que emerge é o de uma \textbf{rede que trabalha em ritmo mais
alto}: as OSS \textbf{atendem 43,7\% mais pacientes}, com os leitos de
internação \textbf{22,7\% mais ocupados} e a UTI \textbf{19,3\% mais
usada} entre os hospitais que a possuem, e as três diferenças são grandes
o bastante para não haver dúvida sobre a sua existência.\footnote{A
intensidade de uso de UTI é estimada com a família Gama, que acomoda a
cauda de UTIs de uso quase nulo; ver a Seção~\ref{sec:inferenciafinal}.}
É uma diferença de regime entre redes com missões distintas, e não um
efeito demonstrado da conversão de um hospital individual (ver as
Seções~\ref{sec:efeitogestao} e~\ref{sec:inferenciafinal}).

O \textbf{paciente fica 10,2\% menos tempo internado} nas OSS quando se
compara cada hospital consigo mesmo, e é isso que fecha a conta do ritmo:
\textbf{mais gente passa pelos mesmos leitos porque cada um fica menos
tempo neles}. A margem de erro sustenta a diferença, mas com cinco
hospitais convertidos ela tende a ser otimista, e o teste mais exigente da
Seção~\ref{sec:inferenciafinal} não a confirma nem a contradiz. Daí a
classificação de sugestiva.

O \textbf{faturamento real por saída} aparece 8,3\% menor, e aqui a
resposta honesta é que \textbf{os dados não decidem}: a margem vai de
23,8\% de queda a 10,4\% de alta, ou seja, comporta hospital mais barato e
hospital mais caro. Com cinco conversões, não há como cravar se as OSS
custam mais ou menos por paciente atendido. \textit{\label{ressalva:pagamento}Ressalva de leitura (regime de pagamento): hospitais OSS são custeados por contrato de gestão com orçamento global, não por reembolso de AIH: o valor registrado por saída é, para eles, produção sem consequência financeira direta; a diferença de faturamento pode refletir o próprio regime de pagamento, a variável cujo efeito se busca medir, e não uso de recursos ou qualidade assistencial. Esta ressalva vale para toda leitura de faturamento no documento e é retomada, em forma abreviada, onde o indicador reaparece.} Os efeitos de ano do mesmo modelo confirmam, agora em
ambiente controlado, a queda real do faturamento da década: 2025 fica 8,9\%
abaixo de 2015 em preços comparáveis.

Na \textbf{mortalidade}, o modelo estima uma queda \textbf{de 5,8\% para
4,6\%} depois da conversão, medida dentro do mesmo hospital, num
patamar de significância que já é limítrofe. \textbf{Esse número não deve
ser citado}: ele não sobrevive ao confronto com o hospital gêmeo da
Seção~\ref{sec:inferenciafinal}. Duas leituras acompanham o resultado e
importam para a gestão. Primeiro, \textbf{as OSS não são hospitais de
mortalidade estruturalmente baixa}: o modelo separa o que é diferença
entre hospitais do que é mudança dentro do hospital, e a parcela entre
hospitais não aponta vantagem alguma; toda a queda estimada vem da
trajetória dos cinco conversores. Segundo, \textbf{nenhum hospital OSS ou
Municipal registra um ano sequer sem óbito} no painel, contra
Filantrópicos e Privados que registram: são redes que recebem casos de
outra natureza, e é mais uma razão para não ler mortalidade bruta como
placar de gestão. Na \textbf{fração de alta complexidade} não há diferença
entre OSS e Direta por método nenhum; quem determina essa fração é a
complexidade instalada do hospital e o seu porte, como era de esperar.

A Tabela~\ref{tab:contrasteoss} traz os mesmos contrastes em número. A
figura seguinte é a leitura visual dela; quem quiser conferir os valores
que sustentam cada afirmação desta seção encontra-os aqui.

\begin{table}[H]
\centering
\caption{O que muda quando o hospital é gerido por OSS em vez da
administração Direta, já descontados complexidade, porte e ano. As linhas
estão ordenadas pela \textbf{firmeza da evidência}, e não pelo tamanho do
número: é a ordem em que um gestor deve confiar neles. A coluna
\textbf{Diferença} é a única que precisa ser lida para decidir; as três
últimas são o registro técnico que sustenta cada linha e estão explicadas
logo abaixo da tabela.}
\label{tab:contrasteoss}
\small
\setlength{\tabcolsep}{4.5pt}
\begin{tabular}{llrrrr}
\toprule
Indicador & Diferença & Evidência & Coef. & EP & $z$ \\
\midrule
\multicolumn{6}{l}{\emph{Diferença firme, medida entre as redes}}\\
Produção (saídas)            & $+43{,}7\%$ de pacientes atendidos & firme & $0{,}3629$ & 0,0895 & $4{,}06$ \\
Ocupação de internação       & $+22{,}7\%$ de leito ocupado & firme & $0{,}2044$ & 0,0779 & $2{,}62$ \\
Ocupação de UTI              & $+19{,}3\%$ de uso da UTI & firme & $0{,}1769$ & 0,0603 & $2{,}93$ \\
\midrule
\multicolumn{6}{l}{\emph{Diferença sugestiva, medida dentro do mesmo hospital}}\\
Tempo de permanência         & $-10{,}2\%$ de tempo internado & sugestiva & $-0{,}1080$ & 0,0468 & $-2{,}31$ \\
\midrule
\multicolumn{6}{l}{\emph{Sem diferença que os dados sustentem}}\\
Mortalidade (geral)          & de 5,8\% para 4,6\%, não sustentada & frágil & $-0{,}2237$ & 0,1311 & $-1{,}71$ \\
Faturamento real por saída   & $-8{,}3\%$ por paciente, dentro da margem & nenhuma & $-0{,}0866$ & 0,0948 & $-0{,}91$ \\
Fração alta complexidade     & de 8,9\% para 7,6\%, dentro da margem & nenhuma & $-0{,}2044$ & 0,2348 & $-0{,}87$ \\
\bottomrule
\end{tabular}
\end{table}

\noindent\textbf{Como ler a coluna Evidência.} \emph{Firme} quer dizer que
a diferença aparece em todos os métodos aplicados e não depende de nenhum
hospital em particular; é retrato de como as redes operam hoje.
\emph{Sugestiva} quer dizer que a diferença é medida comparando cada
hospital consigo mesmo, antes e depois da conversão, e que os cinco
conversores disponíveis bastam para apontá-la mas não para prová-la
(Seções~\ref{sec:efeitogestao} e~\ref{sec:inferenciafinal}).
\emph{Frágil} quer dizer que o número existe no modelo mas some quando o
hospital convertido é comparado a um gêmeo construído a partir dos não
convertidos, e por isso não deve ser citado
(Seção~\ref{sec:inferenciafinal}). \emph{Nenhuma} quer dizer que a
diferença medida é menor que a própria margem de erro: os dados não
distinguem as duas redes nesse indicador.

\comofeito{cada linha é um dos sete modelos principais e as colunas trazem
o contraste entre a gestão por OSS e a administração Direta.
\textbf{Diferença} traduz o resultado do modelo para a unidade que o
gestor usa. Nos indicadores de quantidade ela é a variação percentual.
Nos dois que já são percentuais, a mortalidade e a fração de alta
complexidade, o modelo estima o quanto a taxa se desloca, e a coluna traz
esse deslocamento \textbf{escrito como as duas taxas que ele liga}: a
mortalidade sai de 5,8\% e chega a 4,6\%, a fração de alta complexidade sai
de 8,9\% e chega a 7,6\%. A taxa de partida é a média observada da
administração Direta, que é a rede de referência da comparação. Escrever o
deslocamento como se fosse variação percentual seria erro grosseiro: cair
de 5,8\% para 4,6\% é um deslocamento de 1,2 na escala da taxa, mas
equivale a perder um quinto do valor. \textbf{Coef.} é o valor estimado na escala interna de cada
modelo (o \emph{logito} para as duas proporções, o logaritmo para os
demais), nenhuma delas legível diretamente, e por isso a coluna
\textbf{Diferença} existe. \textbf{EP} é a margem de erro calculada
agrupando as observações por hospital, reconhecendo que 11 anos do mesmo
CNES não são 11 informações independentes. \textbf{$z$} é o coeficiente
dividido pela margem de erro, e valores acima de 2 em módulo indicam
diferença estatisticamente firme. Nas duas proporções, $e^{coef}-1$ seria
razão de chances e não variação percentual da média, e é por isso que ali a
coluna Diferença traz o par de taxas em vez de um único número. A linha da
ocupação de UTI já traz a reestimação pela família Gama da
Seção~\ref{sec:inferenciafinal}.}
\newpage
\comoler{cada linha é um indicador; o quadrado verde é a melhor
estimativa do efeito de ser OSS em vez de Direta, e a haste horizontal
é a margem de erro. Quando a haste cruza a linha vertical do zero, como
no faturamento, não se pode afirmar diferença; quando fica inteira de um
lado, como na produção e na ocupação, a diferença é estatisticamente
firme. Valores à direita do zero indicam OSS maior; à esquerda, OSS
menor.}

\begin{figure}[H]
\centering
\includegraphics[width=.92\textwidth]{fig_est_03_contraste_oss_direta.png}
\caption{Contraste OSS vs Direta nos modelos principais, em efeito
percentual com intervalo de 95\%; a mortalidade, que tem escala
própria, aparece anotada no título. Faturamento e TMP: identificação
\emph{within} (efeitos fixos, sustentada pelos cinco conversores);
produção e ocupações: contraste predominantemente \emph{between} entre
as redes OSS e Direta, de natureza associativa.}
\label{fig:contrasteoss}
\end{figure}

\comofeito{gráfico de pontos com margem de erro. Cada linha é um
indicador (eixo vertical) e a posição do quadrado no eixo horizontal é o
\textbf{efeito estimado de ser OSS em vez de Direta}, em variação
percentual, já descontados complexidade, porte e ano; a haste horizontal
é o intervalo de confiança de 95\%, com as margens agrupadas por
hospital. Quando a haste cruza a linha vertical do zero, não se pode
afirmar diferença. A mortalidade, por ter escala própria, é anotada no
título em vez de entrar no eixo.}
